\section{Caso práctico integrador}
\textbf{Caso académico ficticio elaborado con fines educativos.} Empresa Peruana ABC S.A.C. presenta los siguientes estados financieros resumidos.

\subsection{Estado de Situación Financiera}
\begin{longtable}{lr}
\toprule
Concepto & Importe \\
\midrule
Efectivo & \soles{120,000.00}\\
Cuentas por cobrar & \soles{240,000.00}\\
Inventarios & \soles{180,000.00}\\
Activo corriente & \soles{540,000.00}\\
Activo no corriente & \soles{660,000.00}\\
\textbf{Total activo} & \soles{1,200,000.00}\\
Pasivo corriente & \soles{300,000.00}\\
Pasivo no corriente & \soles{270,000.00}\\
Pasivo total & \soles{570,000.00}\\
Patrimonio & \soles{630,000.00}\\
\textbf{Total pasivo y patrimonio} & \soles{1,200,000.00}\\
\bottomrule
\end{longtable}

\subsection{Estado de Resultados}
\begin{longtable}{lr}
\toprule
Concepto & Importe \\
\midrule
Ventas netas & \soles{1,500,000.00}\\
Costo de ventas & \soles{975,000.00}\\
Utilidad bruta & \soles{525,000.00}\\
Gastos operativos & \soles{300,000.00}\\
Utilidad operativa & \soles{225,000.00}\\
Gastos financieros & \soles{45,000.00}\\
Utilidad antes de impuestos & \soles{180,000.00}\\
Impuesto a la renta & \soles{53,100.00}\\
\textbf{Utilidad neta} & \soles{126,900.00}\\
\bottomrule
\end{longtable}

\subsection{Ratios calculados}
\begin{longtable}{lll}
\toprule
Categoría & Cálculo & Resultado \\
\midrule
Razón corriente & 540,000.00 / 300,000.00 & 1.80 \\
Prueba ácida & (540,000.00 - 180,000.00) / 300,000.00 & 1.20 \\
Endeudamiento & 570,000.00 / 1,200,000.00 x 100 & 47.50 \char37 \\
Deuda-patrimonio & 570,000.00 / 630,000.00 & 0.90 \\
Margen neto & 126,900.00 / 1,500,000.00 x 100 & 8.46 \char37 \\
ROA simple & 126,900.00 / 1,200,000.00 x 100 & 10.58 \char37 \\
ROE simple & 126,900.00 / 630,000.00 x 100 & 20.14 \char37 \\
\bottomrule
\end{longtable}

El ROA y el ROE se presentan como aproximaciones simples porque se usa el saldo final. Para un análisis más riguroso deben emplearse activo promedio y patrimonio promedio.
